
\chapter{طراحی Database و مدل داده}

\section{اصل طراحی}

NFR-11 استفاده از Relational SQL را مشخص می‌کند. Domain نیز به‌شدت Relation-oriented است؛ بنابراین PostgreSQL انتخاب طبیعی است.

\section{User}

\begin{latin}
\begin{lstlisting}
User
- id
- email
- password_hash (managed by Django)
- display_name
- preferred_ui_language
- is_active
- is_staff
- date_joined
- created_at / updated_at if needed
\end{lstlisting}
\end{latin}

\begin{decisionbox}
Custom User Model قبل از اولین Migration واقعی ساخته می‌شود.
\end{decisionbox}

\section{CreatorProfile}

\begin{latin}
\begin{lstlisting}
CreatorProfile
- id
- user_id -> User (OneToOne)
- bio_fa
- bio_en
- avatar
- is_verified
- created_at
\end{lstlisting}
\end{latin}

\section{Topic}

\begin{latin}
\begin{lstlisting}
Topic
- id
- slug
- name_fa
- name_en
- is_active
\end{lstlisting}
\end{latin}

Slug باید unique باشد.

\section{UserInterest}

\begin{latin}
\begin{lstlisting}
UserInterest
- user_id
- topic_id
- created_at
\end{lstlisting}
\end{latin}

Constraint:
\begin{latin}
\begin{lstlisting}
UNIQUE(user_id, topic_id)
\end{lstlisting}
\end{latin}

\section{Source}

Source کار اصلی است که Idea از آن آمده:

\begin{latin}
\begin{lstlisting}
Source
- id
- type: BOOK | ARTICLE | MOVIE | PODCAST
- title
- original_creator_text or structured relation if later needed
- language
- external_url (optional)
- cover (optional)
- metadata (limited)
\end{lstlisting}
\end{latin}

\begin{warningbox}
Source Creator با InsightInk Creator لزوماً یک چیز نیست. مثلاً یک کاربر InsightInk ممکن است درباره کتاب یک نویسنده خارجی Idea بنویسد.
\end{warningbox}

\section{Idea}

\begin{latin}
\begin{lstlisting}
Idea
- id
- creator_id -> CreatorProfile
- source_id -> Source
- language
- title
- body
- estimated_read_seconds
- status: DRAFT | PUBLISHED | HIDDEN
- published_at
- created_at
- updated_at
\end{lstlisting}
\end{latin}

Relation:
\[
Idea \leftrightarrow Topic
\]

\section{Spark}

\begin{latin}
\begin{lstlisting}
Spark
- user_id
- idea_id
- created_at

UNIQUE(user_id, idea_id)
\end{lstlisting}
\end{latin}

\section{SavedIdea}

\begin{latin}
\begin{lstlisting}
SavedIdea
- user_id
- idea_id
- saved_at
- last_opened_at (optional)

UNIQUE(user_id, idea_id)
\end{lstlisting}
\end{latin}

\section{Comment}

\begin{latin}
\begin{lstlisting}
Comment
- id
- author_id -> User
- idea_id -> Idea
- body
- status: VISIBLE | HIDDEN | DELETED
- created_at
- updated_at
\end{lstlisting}
\end{latin}

Nested Comment در Source requirement ضروری نیست و فعلاً اضافه نمی‌شود.

\section{CreatorFollow}

\begin{latin}
\begin{lstlisting}
CreatorFollow
- user_id
- creator_id
- created_at

UNIQUE(user_id, creator_id)
\end{lstlisting}
\end{latin}

\section{ReadingEvent}

برای Story «چند Idea این ماه خواندم» باید اول تعریف Read را مشخص کنیم.

Model احتمالی:
\begin{latin}
\begin{lstlisting}
ReadingEvent
- id
- user_id
- idea_id
- occurred_at
\end{lstlisting}
\end{latin}

\begin{warningbox}
قبل از تعریف دقیق «read»، این Table را نساز. صرف باز شدن صفحه ممکن است معادل خواندن واقعی نباشد.
\end{warningbox}

\section{Report}

\begin{latin}
\begin{lstlisting}
Report
- id
- reporter_id
- target_type
- target_id
- reason
- details
- status
- created_at
- resolved_at
\end{lstlisting}
\end{latin}

برای Generic target باید بین GenericForeignKey و Tableهای مشخص تصمیم بگیریم. برای پروژه آموزشی، سادگی و Queryability مهم‌تر از Generic abstraction است.

\section{ModerationAction}

\begin{latin}
\begin{lstlisting}
ModerationAction
- id
- moderator_id
- report_id (optional)
- action_type
- target information
- note
- created_at
\end{lstlisting}
\end{latin}

\section{Messaging - later}

\begin{latin}
\begin{lstlisting}
Conversation
- id
- created_at

ConversationParticipant
- conversation_id
- user_id

Message
- id
- conversation_id
- sender_id
- body
- created_at
- read_at
\end{lstlisting}
\end{latin}

برای Scope اولیه Conversation بین Reader و Creator است.

\section{Notification - later}

\begin{latin}
\begin{lstlisting}
Notification
- id
- recipient_id
- type
- payload
- read_at
- created_at

NotificationPreference
- user_id
- notification_type
- enabled
\end{lstlisting}
\end{latin}

\section{Membership/eBook - later}

\begin{latin}
\begin{lstlisting}
Ebook
- id
- source_id
- title
- language
- file_metadata
- access_level

Membership
- id
- user_id
- provider_reference
- status
- plan
- period_start
- period_end
\end{lstlisting}
\end{latin}

Raw card data جایی در Schema ما ندارد.

\section{ERD مفهومی}

\begin{center}
\begin{tikzpicture}[
  node distance=1.2cm and 1.5cm,
  entity/.style={draw=InkSoft, rounded corners=2pt, fill=CodeBg, minimum width=2.7cm, minimum height=0.8cm, align=center},
  rel/.style={-{Latex[length=2mm]}, draw=Ash, thick}
]
\node[entity] (user) {User};
\node[entity, right=of user] (creator) {CreatorProfile};
\node[entity, below=of user] (interest) {UserInterest};
\node[entity, right=of interest] (topic) {Topic};
\node[entity, right=of topic] (idea) {Idea};
\node[entity, right=of idea] (source) {Source};
\node[entity, below=of idea] (spark) {Spark / SavedIdea};
\node[entity, below=of creator] (follow) {CreatorFollow};
\node[entity, below=of spark] (comment) {Comment};

\draw[rel] (user) -- node[above]{1:0..1} (creator);
\draw[rel] (user) -- (interest);
\draw[rel] (interest) -- (topic);
\draw[rel] (topic) -- node[above]{M:N} (idea);
\draw[rel] (idea) -- (source);
\draw[rel] (user) -- (spark);
\draw[rel] (spark) -- (idea);
\draw[rel] (user) -- (follow);
\draw[rel] (follow) -- (creator);
\draw[rel] (user) -- (comment);
\draw[rel] (comment) -- (idea);
\end{tikzpicture}
\end{center}

\section{Index Strategy}

Index روی Fieldهایی که مرتب Filter/Sort می‌شوند:
\begin{itemize}
  \item Idea.status؛
  \item Idea.published\_at؛
  \item Idea.language؛
  \item Topic.slug؛
  \item Source.type؛
  \item Notification recipient/read state later.
\end{itemize}

Many-to-many join tableها و Foreign Keyها معمولاً Index مناسب دارند، ولی باید Query plan را بعداً بررسی کنیم.

\section{Normalization}

اصل:
\begin{itemize}
  \item داده تکراری را بی‌دلیل Duplicate نکن؛
  \item Relation واقعی را ForeignKey کن؛
  \item JSONField فقط برای metadata واقعاً flexible؛
  \item Denormalization فقط بعد از Evidence performance.
\end{itemize}

\section{Migration Rules}

\begin{enumerate}
  \item Model را تغییر بده؛
  \item \tech{makemigrations}؛
  \item فایل Migration را بخوان؛
  \item \tech{migrate}؛
  \item Test؛
  \item Commit Model + Migration با هم.
\end{enumerate}
